\chapter{Limitations and Future Work}
\label{ch:limitations-future-work}

% PURPOSE (per ADR-055 amendment 2026-07-31): this chapter absorbs the two layers
% that ADR-054's scope freeze removed from thesis scope, plus the honest limitation
% set. It replaces the June skeleton's empty "Knowledge Graph" and "Analytical
% Query" chapters. Pre-registered but unevaluated hypotheses are reported here --
% explicitly, never silently dropped (no-HARKing convention).

\section{Limitations}
\label{sec:limitations}

\subsection{Corpus: Synthetic and Born-Digital}
\label{sec:lim-corpus}
% TODO: the evaluation corpus is generated from the reference e-invoicing corpus and
% rendered as born-digital PDFs. Real firm documents arrive as scans and phone
% photographs, with skew, stamps, staples and handwriting. Published multilingual
% parsing benchmarks report a substantial accuracy drop for open-source parsers on
% *photographed* documents, so the numbers reported here are an upper bound for the
% camera-scanned case. State this before any result is interpreted.

\subsection{In-Sample versus Held-Out Measurement}
\label{sec:lim-insample}
% TODO: distinguish the diagnostic (in-sample) runs used to build the measurement
% apparatus from the sealed split used for the reported baseline. Name explicitly
% which numbers in Ch.~\ref{ch:results} are which, and why an in-sample number may
% never be quoted as real-world accuracy.

\subsection{Hardware Envelope}
\label{sec:lim-hardware}
% TODO: local-first inference was developed inside a 16\,GB unified-memory laptop;
% one candidate exceeded that envelope and had to be measured on rented hardware.
% The privacy premise is on-premises processing, which permits a workstation GPU --
% but the laptop-only claim is narrower than "runs locally".

\subsection{Language and Document Class}
\label{sec:lim-scope}
% TODO: German B2B invoices only. No claim of transfer to other languages, other
% document classes (contracts, bank statements, receipts) or other jurisdictions.

\subsection{Measurement Constructs}
\label{sec:lim-measurement}
% TODO: the reading-quality proxy (whether a ground-truth value is findable in the
% transcript) correlates strongly with the end metric but is a proxy: a value can be
% present as a string yet unusable because the surrounding table structure is
% mangled. State the correlation honestly and mark the proxy as a proxy.

\section{Pre-Registered but Unevaluated Hypotheses}
\label{sec:unevaluated}
% TODO (no-HARKing): list the hypotheses that were registered before data collection
% but not tested within thesis scope -- the knowledge-graph fidelity claim, the
% graph-versus-vector retrieval claim, the query-routing claim, the conditional
% validator-retry claim, the full cloud-versus-local comparison, and the
% template-shift claim. Report them as "not evaluated within thesis scope" with the
% reason (scope freeze in favour of a complete, honest Layer-1 result), NOT as
% negative or inconclusive findings.

\section{Future Work}
\label{sec:future-work}

\subsection{Knowledge-Graph Construction}
\label{sec:fw-kg}
% TODO: build Layer 2 per Sec.~\ref{sec:design-layer2}; measure how extraction error
% propagates into graph fidelity (does a field-level error rate predict a
% compliance-decision error rate?). This is the most direct continuation.

\subsection{Analytical Query and Retrieval Comparison}
\label{sec:fw-query}
% TODO: build Layer 3 per Sec.~\ref{sec:design-layer3}; the open question is which
% question classes genuinely benefit from graph structure over flat retrieval.

\subsection{Cloud Comparison Arm}
\label{sec:fw-cloud}
% TODO: an EU-hosted cloud comparison would quantify the accuracy cost of the
% privacy-first constraint. Deliberately excluded here because sending client
% documents to a third party is the exact practice the thesis premise rules out;
% a comparison would need synthetic documents only.

\subsection{Robustness on Degraded Documents}
\label{sec:fw-robustness}
% TODO: scanned/photographed robustness, following the multilingual-parsing
% benchmark's photographed-document protocol.

\subsection{Adaptation to Further Layouts and Languages}
\label{sec:fw-transfer}
% TODO: template-shift and cross-language transfer; low-rank adaptation
% (\ac{LoRA}) of the structurer as a cheap adaptation path once a reliable reader
% is fixed.

\subsection{Production Hardening}
\label{sec:fw-production}
% TODO: throughput, batching, audit logging, and the human-in-the-loop review
% workflow a firm would actually need.
